\chapter{Curvas Planas}\label{ch:b1:curves}

O cálculo do \cref{ch:b1:derivative,ch:b1:taylor} foi construído para funções $y = f(x)$; a maioria das
curvas da geometria e da mecânica --- trajetórias, círculos rolando
sobre círculos, órbitas --- recusa essa forma e apresenta-se antes
como curvas \emph{parametrizadas} $t \mapsto (x(t),
y(t))$ ou como curvas \emph{polares}
$r = r(\theta)$. Este capítulo clássico estuda as duas, e fecha com as \omterm{def:b1:curves:conics}{cônicas}.

\section{Curvas parametrizadas}

\begin{definition}\label{def:b1:curves:def}
Uma \emph{curva parametrizada}\index{curva parametrizada} é uma
\omterm{def:b1:logic:map}{aplicação} $f
\colon I \to \R^2$, $t \mapsto M(t) = (x(t), y(t))$, com $x, y$ de classe $C^1$ (no mínimo) no
\omterm{prop:b1:reals:intervals}{intervalo} $I$. O \emph{vetor velocidade} é $f'(t) = (x'(t), y'(t))$; o ponto $M(t_0)$ é
\emph{regular} quando $f'(t_0) \neq (0,0)$, e a \emph{tangente} nesse ponto é a reta
por $M(t_0)$ dirigida por $f'(t_0)$.
\end{definition}

\begin{proof}[Por que $f'$ dirige a tangente]
Por Taylor--Young componente a componente, $M(t) = M(t_0) + (t - t_0) f'(t_0) +
o(t - t_0)$: a direção da corda
$\frac{M(t) - M(t_0)}{t - t_0}$ tende a $f'(t_0)$. A hipótese de regularidade é o que faz do
limite uma \emph{direção}: se $f'(t_0) = 0$, a fórmula exibida colapsa em
$M(t) = M(t_0) + o(t - t_0)$ e nada diz sobre como a curva deixa o ponto --- daí o
tratamento separado dos pontos singulares no método abaixo, em que a
primeira \omterm{def:b1:derivative:def}{derivada} \emph{não nula} assume o papel de $f'$. Note também
que a reta tangente é um objeto geométrico: toda reparametrização
muda $f'$ por um fator escalar não nulo e deixa a reta inalterada.
\end{proof}

\begin{method}[Estudar uma curva parametrizada]\label{met:b1:curves:study}
\begin{enumerate}
  \item \emph{Reduza o domínio} usando simetrias: relações entre
        $M(-t)$, $M(t + T)$, \dots\ e $M(t)$ (reflexões, translações,
        periodicidade), e depois desenhe apenas a parte reduzida.
  \item \emph{Variações}: tabule os sinais de $x'$ e de $y'$
        conjuntamente; a curva move-se para a direita ou para a
        esquerda conforme $x'$, e para cima ou para baixo conforme
        $y'$.
  \item \emph{Pontos notáveis}: tangente horizontal ($y' = 0 \neq
        x'$),
        tangente vertical ($x' = 0 \neq y'$); num ponto \emph{singular}
        ($f' = 0$), examine as \omterm{def:b1:derivative:def}{derivadas} de ordem superior para
        obter a direção da tangente.
  \item \emph{Comportamento assintótico} nas extremidades de $I$, e
        depois esboce.
\end{enumerate}
\end{method}

\begin{example}[Ramos assintóticos, trabalhados]
\label{ex:b1:curves:asymptotes}
$x(t) = t$, $y(t) = t + \dfrac1t$ em $\intoo{0}{+\infty}$. Quando $t
\to 0^{+}$: $x \to 0$ enquanto $y \to +\infty$ --- a
curva sobe ao longo da \emph{assíntota vertical} $x = 0$ (limite
finito para uma coordenada, infinito para a outra). Quando $t \to +\infty$: as
duas \omterm{prop:b1:vspaces:coordinates}{coordenadas} explodem, logo teste uma reta: a diferença
\[
y(t) - x(t) = \frac1t \longrightarrow 0^{+}
\]
exibe a \emph{assíntota oblíqua} $y = x$, aproximada por cima. No
meio disso, $y' = 1 - \frac1{t^2}$ anula-se em $t = 1$: o ponto $(1, 2)$ é o ponto
mais baixo do ramo, e $x' = 1 >
0$ em toda parte, logo a curva avança sempre
para a direita. Um arco, duas assíntotas, um mínimo: um quadro
completo a partir de três cálculos --- com a receita geral visível
por baixo: quando $x, y \to \infty$ juntos, examine $y/x$ à procura de um
coeficiente angular candidato (aqui $\to 1$), e depois $y - (\text{inclinação})\,x$ para a
ordenada na origem e para o lado da aproximação.
\end{example}

\begin{example}[A astroide]\label{ex:b1:curves:astroid}
$x(t) = \cos^3 t$, $y(t) = \sin^3 t$.\index{astroide} Simetrias: $M(t + 2\pi) = M(t)$ (estude um período);
$M(-t)$ é o reflexo de $M(t)$ no eixo $x$; $M(\pi - t)$, no eixo $y$; $M(\frac\pi2 -
t)$,
na diagonal $y = x$: basta estudar $t \in
\intcc{0}{\frac\pi4}$ e desdobrar.

Velocidade: $f'(t) = 3\sin t\cos t\,(-\cos t, \sin t)$. Em $\intoo{0}{\frac\pi2}$ todos os pontos são regulares, com
tangente dirigida por $(-\cos t, \sin t)$; em $t = 0$ (o ponto $(1,0)$) a
velocidade se anula: uma \emph{cúspide}\index{cúspide}, em que a
curva inverte o sentido ao longo da direção tangente $(-1, 0)\cdot(\pm)$ --- por
simetria, as quatro \omterm{ex:b1:curves:astroid}{cúspides} ficam em $(\pm1, 0), (0, \pm1)$.
\end{example}

\begin{example}[Um oito, estudado por completo]
\label{ex:b1:curves:lissajous}
$x(t) = \sin t$, $y(t) = \sin 2t$ (uma \emph{curva de Lissajous}). \emph{Simetrias}:
$M(t + \pi) = (-x(t), y(t))$ (reflexão no eixo $y$), $M(-t) = (-x(t), -y(t))$ (simetria central), $M(\pi - t) = (x(t), -y(t))$
(reflexão no eixo $x$): basta estudar $t \in \intcc{0}{\frac\pi2}$ e desdobrar.
\emph{Variações}: $x' = \cos t \geq 0$ em toda parte, ao passo que $y' =
2\cos 2t$ é positiva
antes de $t = \frac\pi4$ e negativa depois: o arco sobe para a direita até o
cume $\bigl(\frac{\sqrt2}2,\ 1\bigr)$ em $t = \frac\pi4$ (tangente horizontal), e depois desce para a
direita até $(1, 0)$ em $t =
\frac\pi2$, onde $x' = 0 \neq y'$: tangente vertical.
\emph{Ponto duplo}: $M(0) = M(\pi) = (0,0)$, com duas velocidades diferentes
\[
f'(0) = (1,\ 2), \qquad f'(\pi) = (-1,\ 2) :
\]
dois ramos \emph{regulares} cruzando-se na origem sob ângulos
distintos --- um ponto duplo, não um ponto singular: cada passagem é
perfeitamente suave, e as duas passagens apenas partilham a sua
localização. A curva inteira é o oito abaixo.
\end{example}

\begin{omfigure}
\begin{tikzpicture}[scale=1.5]
  \draw[->, thin] (-1.25,0) -- (1.3,0) node[below] {$x$};
  \draw[->, thin] (0,-1.25) -- (0,1.3) node[left] {$y$};
  \draw[omDef, very thick, domain=0:360, samples=200, smooth]
    plot ({sin(\x)}, {sin(2*\x)});
  \fill (0.7071, 1) circle (0.035);
  \node[above, font=\small] at (0.7071, 1)
    {$\bigl(\tfrac{\sqrt2}2, 1\bigr)$};
  \fill (0,0) circle (0.035);
\end{tikzpicture}

{\small A curva de Lissajous $(\sin t, \sin 2t)$: um ponto duplo na origem, onde
dois ramos regulares se cruzam com velocidades $(1, 2)$ e $(-1, 2)$,
e tangentes horizontais nos quatro cumes. A simetria reduziu todo o
trabalho a um quarto de período.}
\end{omfigure}

\begin{omfigure}
\begin{tikzpicture}[scale=1.35]
  \draw[->, thin] (-1.3,0) -- (1.35,0) node[below] {$x$};
  \draw[->, thin] (0,-1.3) -- (0,1.35) node[left] {$y$};
  \draw[omDef, very thick, domain=0:360, samples=200, smooth]
    plot ({cos(\x)^3}, {sin(\x)^3});
  \node[below right, font=\small] at (1,0) {$1$};
  \node[above left, font=\small] at (0,1) {$1$};
\end{tikzpicture}

{\small A astroide $(\cos^3 t, \sin^3 t)$: quatro arcos que se encontram em quatro
\omterm{ex:b1:curves:astroid}{cúspides}. É a curva traçada por um ponto de um círculo de raio
$\frac14$ rolando dentro do círculo unitário.}
\end{omfigure}

\section{Curvas polares}

\begin{definition}\label{def:b1:curves:polar}
Uma \emph{curva polar}\index{curva polar} é dada por $r = r(\theta)$: o ponto
de parâmetro $\theta$ é
\[
M(\theta) = r(\theta)\,\vec u(\theta),
\qquad \vec u(\theta) = (\cos\theta, \sin\theta).
\]
Com $\vec v(\theta) = (-\sin\theta, \cos\theta) = \vec
u\,'(\theta)$, a velocidade é
\[
M'(\theta) = r'(\theta)\, \vec u(\theta) + r(\theta)\, \vec
v(\theta) .
\]
Consequências: onde $r \neq 0$, o ponto é regular, e a tangente faz com o
raio o ângulo $V$ dado por $\tan V =
\frac{r}{r'}$ (ângulo entre $M'$ e $\vec u$); onde
$r(\theta_0)
= 0$, a curva passa pela origem \emph{com tangente o raio $\theta = \theta_0$}
(direção $\vec u(\theta_0)$, lida em $M'(\theta_0) = r'(\theta_0)\vec u(\theta_0)$, ou a partir da corda limite: a corda
de $O$ a $M(\theta)$ é carregada pelo próprio $\vec
u(\theta)$, que tende a $\vec u(\theta_0)$ ---
de modo que a regra vale mesmo quando $r'(\theta_0) = 0$ e o ponto é singular,
razão pela qual as \omterm{ex:b1:curves:astroid}{cúspides} polares na origem, como a da cardioide,
ganham a sua tangente de graça).
\end{definition}

\begin{example}[A cardioide]\label{ex:b1:curves:cardioid}
$r(\theta) = 1 + \cos\theta$.\index{cardioide} Simetria: $r(-\theta)
= r(\theta)$: reflexão no eixo $x$; estude
$\theta \in
\intcc{0}{\pi}$. $r$ decresce de $2$ a $0$; em $\theta = \pi$, $r = 0$: a curva atinge a
origem tangencialmente ao raio $\theta = \pi$ (o eixo $x$), formando ali uma
\omterm{ex:b1:curves:astroid}{cúspide} --- a ponta do coração. Tangente em $\theta = 0$: $r' = 0$, logo
$\tan V = \infty$: perpendicular ao eixo.

O ângulo $V$ merece mais uma leitura. Em $\theta =
\frac\pi2$: $r = 1$ e $r' = -1$,
logo $\tan V = \frac{r}{r'} =
-1$: a tangente faz três quartos de um ângulo reto com o raio
que sai --- a curva já está dobrando de volta rumo à sua \omterm{ex:b1:curves:astroid}{cúspide}. A
fórmula $\tan V = r/r'$ entrega direções tangentes ao longo de toda a curva sem
calcular $M'(\theta)$ coisa alguma: é o análogo polar de ler um coeficiente
angular.
\end{example}

\begin{example}[De polar a cartesiana: um círculo escondido]
\label{ex:b1:curves:polarcircle}
O que é a \omterm{def:b1:curves:polar}{curva polar} $r = 2\cos\theta$? Multiplique por $r$: $r^2 = 2r\cos\theta$, isto é,
$x^2 + y^2 = 2x$, isto é,
\[
(x - 1)^2 + y^2 = 1 :
\]
o círculo de centro $(1, 0)$ e raio $1$, que passa pela origem. A
contabilidade importa: quando $\theta$ percorre $\intoc{-\frac\pi2}{\frac\pi2}$, o círculo
inteiro é varrido exatamente uma vez ($r$ anula-se nas duas
extremidades), e em $\theta =
\pm\frac\pi2$ a regra ``tangente na origem ao longo do
raio $\theta = \theta_0$'' dá ali uma tangente \emph{vertical} --- coerente com a
geometria, já que o eixo vertical é de fato tangente a este círculo
em $O$. Para $\theta$ além desse \omterm{prop:b1:reals:intervals}{intervalo}, $r < 0$ percorre de novo
o mesmo círculo: um lembrete de que uma \omterm{def:b1:curves:polar}{curva polar} é um objeto
\emph{parametrizado}, autorizado a passar por cima de si mesmo.
\end{example}

\begin{omfigure}
\begin{tikzpicture}[scale=1.5]
  \draw[->, thin] (-1.2,0) -- (1.3,0) node[below] {$x$};
  \draw[->, thin] (0,-1.2) -- (0,1.3) node[left] {$y$};
  \draw[omProp, very thick, domain=0:360, samples=300, smooth]
    plot ({cos(2*\x)*cos(\x)}, {cos(2*\x)*sin(\x)});
  \draw[thin, dashed] (-0.85,-0.85) -- (0.85,0.85);
  \draw[thin, dashed] (-0.85,0.85) -- (0.85,-0.85);
\end{tikzpicture}

{\small A rosácea de quatro pétalas $r = \cos 2\theta$ (o \cref{exo:b1:curves:4}). As pétalas ao
longo do eixo $y$ são traçadas com $r < 0$ (o ponto é marcado no raio
oposto a $\theta$); as diagonais tracejadas $\theta = \pm\frac\pi4$ são as tangentes na
origem, onde $r$ se anula.}
\end{omfigure}

\begin{omfigure}
\begin{tikzpicture}[scale=1.15]
  \draw[->, thin] (-0.7,0) -- (2.4,0) node[below] {$x$};
  \draw[->, thin] (0,-1.7) -- (0,1.7) node[left] {$y$};
  \draw[omThm, very thick, domain=0:360, samples=200, smooth]
    plot ({(1+cos(\x))*cos(\x)}, {(1+cos(\x))*sin(\x)});
  \node[below, font=\small] at (2,0) {$2$};
\end{tikzpicture}

{\small A cardioide $r = 1 + \cos\theta$. As \omterm{def:b1:curves:polar}{curvas polares} leem-se varrendo o
ângulo: o raio incha e encolhe à medida que $\theta$ gira.}
\end{omfigure}

\section{Comprimento de arco}

\begin{definition}[Comprimento de um arco]\label{def:b1:curves:arclength}
O \emph{comprimento}\index{comprimento de arco} de um arco $C^1$
$f \colon
\intcc{a}{b} \to \R^2$ é a \omterm{thm:b1:integration:def}{integral} da velocidade escalar:
\[
L = \int_a^b \norm{f'(t)}\,\dd t
= \int_a^b \sqrt{x'(t)^2 + y'(t)^2}\;\dd t .
\]
(Motivação: num \omterm{prop:b1:reals:intervals}{intervalo} pequeno, $M(t + h) \approx M(t) + h
f'(t)$, logo o arco é próximo de
uma poligonal cujos comprimentos de segmento somam uma soma de
Riemann de $\norm{f'}$, o \cref{thm:b1:integration:riemann}.) Para uma \omterm{def:b1:curves:polar}{curva polar} $r =
r(\theta)$, a velocidade
$r'\vec u + r\vec v$ tem componentes \omterm{def:b1:euclid:orthogonal}{ortogonais}, logo
\[
L = \int_{\theta_1}^{\theta_2}
\sqrt{r'(\theta)^2 + r(\theta)^2}\;\dd\theta .
\]
O comprimento não depende da parametrização (monótona, $C^1$)
escolhida: substituir $t = \varphi(s)$ na \omterm{thm:b1:integration:def}{integral} (o \cref{thm:b1:integration:parts}) multiplica $f'$ por
$\varphi'$ e $\dd t$ por $\varphi'^{-1}$.
\end{definition}

\begin{example}[Verificação de coerência: o círculo]\label{ex:b1:curves:circlelength}
Para $f(t) = (R\cos t, R\sin t)$ em $\intcc{0}{2\pi}$: $\norm{f'} =
R$, logo $L = 2\pi R$ --- a definição devolve a
circunferência. Teste de reparametrização: a \omterm{def:b1:logic:map}{aplicação} $g(t) = (R\cos 2t, R\sin 2t)$ em $\intcc{0}{\pi}$
desenha o \emph{mesmo} círculo ao dobro da velocidade $\norm{g'} = 2R$, e
\[
\int_0^{\pi} 2R\,\dd t = 2\pi R
\]
de novo: metade do tempo, o dobro da velocidade, o mesmo \omterm{def:b1:curves:arclength}{comprimento}
--- a invariância prometida na definição, observada uma vez em
números. (Rodar $g$ em todo o $\intcc{0}{2\pi}$ daria $4\pi R$: uma curva
percorrida duas vezes é duas vezes mais longa \emph{como percurso};
o \omterm{def:b1:curves:arclength}{comprimento} mede o caminho parametrizado, e a contabilidade honesta
do \omterm{prop:b1:reals:intervals}{intervalo} faz parte do cálculo.) Para a astroide $(\cos^3t, \sin^3t)$: $\norm{f'} = 3\abs{\sin
t\cos t} = \tfrac32\abs{\sin 2t}$, e
por simetria $L = 4\int_0^{\pi/2}
\tfrac32\sin 2t\,\dd t = 6$: uma curva desenhada dentro do círculo unitário,
de \omterm{def:b1:curves:arclength}{comprimento} $6 < 2\pi$. O problema de fim de semana mede o arco mais
famoso de todos.
\end{example}

\begin{omfigure}
\begin{tikzpicture}[scale=0.85]
  \draw[->, thin] (-0.5,0) -- (7.2,0) node[below] {$x$};
  \draw[->, thin] (0,-0.4) -- (0,2.6) node[left] {$y$};
  \draw[omThm, very thick, domain=0:6.2832, samples=160, smooth]
    plot ({\x - sin(\x r)}, {1 - cos(\x r)});
  \draw[omDef, thick] (2,1) circle (1);
  \fill[omDef] (2,1) circle (0.04);
  \fill (2 - 0.9093, 1.4161) circle (0.06);
  \draw[omProp, thick] (2,0) -- (2 - 0.9093, 1.4161) -- (2,2);
  \fill (2,0) circle (0.05);
  \fill (2,2) circle (0.05);
  \node[below, font=\small] at (2,0) {$C$};
  \node[above, font=\small] at (2,2) {$T$};
  \node[left, font=\small] at (2 - 0.9093, 1.4161) {$M$};
  \node[below, font=\small] at (6.2832,0) {$2\pi R$};
\end{tikzpicture}

{\small Um arco da cicloide (raio $R = 1$), o círculo rolante em
$t = 2$, e as duas retas guia do problema de fim de semana: a corda
$MC$ até o ponto de contato é \emph{normal} à curva, e a corda $MT$
até o topo do círculo é \emph{tangente}.}
\end{omfigure}

\section{Cônicas}

\begin{definition}[Definição por foco e diretriz]\label{def:b1:curves:conics}
Fixe um ponto $F$ (\emph{foco}), uma reta $D$ que não passa por $F$
(\emph{diretriz}) e $e > 0$ (\emph{excentricidade}). A
\emph{cônica}\index{cônica} destes dados é
\[
\mathcal{C} = \{M : d(M, F) = e\; d(M, D)\}:
\]
uma \emph{elipse} para $e < 1$, uma \emph{parábola} para $e = 1$, uma
\emph{hipérbole} para $e > 1$. (O círculo aparece como o limite
degenerado $e \to 0$.)
\end{definition}

\begin{example}[A definição, conferida numa parábola]
\label{ex:b1:curves:parabolacheck}
Tome a parábola $y^2 = 4x$, isto é, $2p = 4$: foco $F = (1, 0)$ e diretriz $D : x = -1$
(a forma reduzida do \cref{thm:b1:curves:reduced} os coloca em $\pm\frac p2$). No ponto $M = (1, 2)$ da
curva:
\[
MF = \sqrt{(1-1)^2 + 2^2} = 2,
\qquad
d(M, D) = 1 - (-1) = 2 :
\]
iguais, como $e = 1$ exige. Em $M' = (4, 4)$: $MF' = \sqrt{9 +
16} = 5$ e $d(M', D) = 5$ de novo. A
definição por foco e diretriz não é uma abstração: é um par de
distâncias que se pode medir em qualquer ponto, e a álgebra das
equações reduzidas nada mais é do que essa medição feita uma vez por
todas.
\end{example}

\begin{theorem}[Equações reduzidas]\label{thm:b1:curves:reduced}
Num referencial \omterm{def:b1:euclid:orthogonal}{ortonormal} bem escolhido:
\[
\text{elipse: } \frac{x^2}{a^2} + \frac{y^2}{b^2} = 1
\quad (a \geq b > 0),
\qquad
\text{hipérbole: } \frac{x^2}{a^2} - \frac{y^2}{b^2} = 1,
\qquad
\text{parábola: } y^2 = 2px,
\]
com, para a elipse: focos em $(\pm c, 0)$, $c = \sqrt{a^2 - b^2}$, $e = \frac ca$, e a caracterização
bifocal $MF + MF' = 2a$; para a hipérbole: $c = \sqrt{a^2 + b^2}$, $e = \frac ca$, $\abs{MF - MF'} = 2a$, assíntotas
$y = \pm\frac ba x$.
\end{theorem}

\begin{proof}
Tome o foco na origem e a diretriz vertical, $x = -h$ ($h > 0$): a
condição $MF^2 = e^2 d(M, D)^2$ lê-se $x^2 + y^2 =
e^2 (x + h)^2$. Para $e = 1$: $y^2 = 2hx + h^2$, uma parábola depois
do deslocamento $x \mapsto x - \frac h2$ (logo $p = h$). Para $e \neq 1$: completando o
quadrado em $x$,
\[
(1 - e^2)\Bigl(x - \frac{e^2 h}{1 - e^2}\Bigr)^{\!2} + y^2
= e^2h^2 + \frac{e^4 h^2}{1 - e^2}
= \frac{e^2 h^2}{1 - e^2} .
\]
Ponha $X = x - \frac{e^2h}{1-e^2}$ (um deslocamento do referencial). Quando $e < 1$, divida
pelo lado direito, que é positivo: $\frac{X^2}{a^2} + \frac{y^2}{b^2} = 1$ com $a = \frac{eh}{1 -
e^2}$, $b = \frac{eh}{\sqrt{1-e^2}}$; quando
$e > 1$, os dois lados da fórmula exibida trocam de sinal
adequadamente e a mesma divisão dá $\frac{X^2}{a^2} - \frac{y^2}{b^2} = 1$ com $a = \frac{eh}{e^2 -
1}$, $b = \frac{eh}{\sqrt{e^2 - 1}}$. Os valores
indicados de $c$ seguem ($c^2 = a^2 - b^2$ ou $a^2 + b^2$ dá $c = ea$ nos dois casos,
colocando o foco corretamente), e as propriedades bifocais são
verificações diretas nas equações reduzidas. Em detalhe para a
elipse, com $F' = (c, 0)$ e $M = (X, y)$ sobre a curva:
\[
MF'^2 = (X - c)^2 + y^2
= X^2 - 2cX + c^2 + b^2\Bigl(1 - \frac{X^2}{a^2}\Bigr)
= \frac{c^2}{a^2}X^2 - 2cX + a^2
= (a - eX)^2 ,
\]
usando $c^2 + b^2 = a^2$ e $e = \frac ca$; como $\abs X \leq
a$ e $e < 1$, $a - eX > 0$, logo $MF' = a - eX$, sem
nunca extrair raiz quadrada alguma. O cálculo espelhado dá $MF = a + eX$,
donde $MF + MF' = 2a$, constante. Para a hipérbole a mesma álgebra dá $MF = \abs{a + eX}$ e
$MF' = \abs{eX - a}$, com diferença $\pm2a$ conforme o ramo.
\end{proof}

\begin{example}\label{ex:b1:curves:conicexample}
$\frac{x^2}{25} + \frac{y^2}{9} = 1$: elipse, $a = 5$, $b = 3$, $c = 4$: focos $(\pm4, 0)$, excentricidade
$\frac45$. A sua parametrização: $(5\cos t, 3\sin t)$ --- um círculo esticado
anisotropicamente; a soma das distâncias aos focos, para qualquer um
dos seus pontos, é $10$.
\end{example}

\begin{example}[Ler uma órbita a partir da sua equação polar]
\label{ex:b1:curves:orbit}
A \omterm{def:b1:curves:conics}{cônica} polar $r = \dfrac{1}{1 + \frac12\cos\theta}$ (o \cref{exo:b1:curves:7} com $p = 1$, $e = \frac12$) é uma elipse com um
\emph{foco na origem} --- a geometria de uma órbita planetária com o
sol em $O$. Extraia tudo de $p$ e de $e$: pela redução do \cref{thm:b1:curves:reduced},
$a = \dfrac{p}{1 - e^2} =
\dfrac{1}{3/4} = \dfrac43$ e $c = ea = \dfrac23$. As duas ápsides conferem isso sem teoria alguma:
\[
r(0) = \frac{1}{3/2} = \frac23 = a - c
\quad (\text{periélio}),
\qquad
r(\pi) = \frac{1}{1/2} = 2 = a + c
\quad (\text{afélio}),
\]
e a sua soma $\frac23 + 2 = \frac83 = 2a$ recupera o eixo maior. A forma polar é a natural
sempre que um foco é fisicamente distinguido; a forma cartesiana
reduzida, sempre que o são os eixos de simetria. Converter entre as
duas é exatamente o que faz o cálculo do quadrado completado do
teorema.

A excentricidade como um botão, por fim: mantenha $p = 1$ e gire $e$
em $r
= \frac{1}{1 + e\cos\theta}$. Em $e = 0$: o círculo $r = 1$. Em $e = \frac12$: a elipse que
acabamos de estudar, com $r$ oscilando entre $\frac23$ e $2$. Em
$e = 1$: $r(\theta) \to \infty$ quando $\theta
\to \pi$ --- a curva já não se fecha: uma
parábola, com o seu ponto mais distante empurrado para o infinito. Em
$e = 2$: o denominador anula-se em $\cos\theta = -\frac12$, e só $\theta \in
\intoo{-\frac{2\pi}3}{\frac{2\pi}3}$ sobrevive: um
ramo de hipérbole, escapando ao longo de duas direções assintóticas.
Uma fórmula, a família das \omterm{def:b1:curves:conics}{cônicas} inteira, e o ponto de transição
$e = 1$ visível como o instante em que o denominador atinge zero pela
primeira vez.
\end{example}

\begin{remark}[Armadilhas comuns]
\label{rem:b1:curves:pitfalls}
\emph{Pontos duplos não são pontos singulares}: num autocruzamento,
cada ramo é regular; ``singular'' refere-se a $f'(t_0) = 0$ para um valor do
parâmetro (o \cref{ex:b1:curves:lissajous}, contra as \omterm{ex:b1:curves:astroid}{cúspides} da astroide).
\emph{Tangente vertical versus \omterm{ex:b1:curves:astroid}{cúspide}}: $x'(t_0) = 0 \neq y'(t_0)$ é um ponto regular com
tangente vertical; só $x' = y' = 0$ exige a expansão de ordem superior.
\emph{Raios negativos marcam-se no raio oposto}: para $r(\theta) <
0$ o ponto é
$-\abs{r}\,\vec u(\theta)$, no ângulo $\theta +
\pi$ (o \cref{exo:b1:curves:5}); esquecer isso perde laços internos
ou duplica curvas.
\emph{O \omterm{def:b1:curves:arclength}{comprimento de arco} integra $\norm{f'}$, e não $f'$}: separe a
\omterm{thm:b1:integration:def}{integral} nos zeros da velocidade escalar --- para a astroide,
integrar $\frac32\sin 2t$ ao longo de um período completo sem valores
absolutos dá $0$, e não $6$.
\emph{As \omterm{def:b1:curves:conics}{cônicas} têm de ser reduzidas antes de serem lidas}: em
$9x^2 +
25y^2 - 36x - 50y - 164 = 0$, nem os eixos, nem o centro, nem a excentricidade são visíveis
antes de completar os quadrados (o \cref{exo:b1:curves:6}); e uma parte quadrática
que se anula numa variável significa parábola, não ``elipse
degenerada''.
\end{remark}

\begin{remark}[Para onde vão estas curvas]
\label{rem:b1:curves:whereused}
As \omterm{def:b1:curves:def}{curvas parametrizadas} são a língua da mecânica: as trajetórias são
curvas, os vetores velocidade são vetores tangentes, e o \omterm{def:b1:curves:arclength}{comprimento
de arco} da \cref{def:b1:curves:arclength} é a distância percorrida. As \omterm{def:b1:curves:conics}{cônicas} reaparecem onde
quer que atue uma lei do inverso do quadrado --- as órbitas
planetárias são elipses com o sol num foco. No \cref{ch:b1:multivar}, as curvas
tornam-se os \omterm{def:b1:logic:sets}{conjuntos} de nível de funções de duas variáveis, e a
tangente deste capítulo encontra o gradiente do próximo. O volume do
segundo ano de graduação acrescenta a curvatura e a forma canônica
local de uma curva; o problema de fim de semana abaixo já extrai,
apenas com as ferramentas deste ano, tudo o que o século XVII sabia
sobre a sua curva mais celebrada.
\end{remark}

\begin{remark}[Perspectivas dentro do Livro 3]
\label{rem:b1:curves:perspectives}
Este capítulo consome toda a primeira metade do volume e alimenta o
último capítulo. \emph{Consumido}: os vetores tangentes são \omterm{def:b1:derivative:def}{derivadas}
(o \cref{ch:b1:derivative}), o \omterm{def:b1:curves:arclength}{comprimento de arco} e as áreas são integrais (o \cref{ch:b1:integration}),
as \omterm{ex:b1:curves:astroid}{cúspides} são resolvidas por expansões de Taylor (o \cref{ch:b1:taylor}), a
tautócrona é uma equação diferencial \omterm{def:b1:linmaps:def}{linear} (o \cref{ch:b1:diffeq}), e toda
distância e todo ângulo são euclidianos (o \cref{ch:b1:euclid}).
\emph{Alimentado}: no \cref{ch:b1:multivar}, uma curva $t \mapsto (x(t),
y(t))$ traçada dentro de um
\omterm{def:b1:logic:sets}{conjunto} de nível $\{f = c\}$ é \omterm{def:b1:derivative:def}{derivada} pela regra da cadeia, e a
identidade resultante $\langle \nabla f,\
f'(t)\rangle = 0$ casa os vetores tangentes deste capítulo
com os gradientes do próximo --- a reta tangente da curva e a direção
normal da superfície são o mesmo cálculo visto das duas margens.
\end{remark}

\section{Exercícios}

\begin{exercise}[$\star$]\label{exo:b1:curves:1}
Estude e esboce a curva $x(t) = t^2$, $y(t) = t^3$ (simetrias, variações,
comportamento no ponto singular $t = 0$).
\end{exercise}

\begin{exercise}[$\star$]\label{exo:b1:curves:2}
Para a cicloide $x(t) = t - \sin t$, $y(t) = 1 - \cos t$ (a trajetória de um ponto de uma roda de
raio $1$ que rola): identifique a simetria de translação por período,
os pontos singulares, e a direção da tangente em $t = 0$
\emph{(expanda $x$ e $y$ até as primeiras ordens não nulas)}.
\end{exercise}

\begin{exercise}[$\star$]\label{exo:b1:curves:3}
Dê a reta tangente à curva $(\cos^3 t, \sin^3 t)$ em $t =
\frac\pi4$, e verifique que o segmento
dessa tangente cortado pelos eixos tem \omterm{def:b1:curves:arclength}{comprimento} $1$ --- uma
propriedade famosa da astroide (verdadeira em todo ponto regular).
\end{exercise}

\begin{exercise}[$\star$]\label{exo:b1:curves:4}
Esboce as \omterm{def:b1:curves:polar}{curvas polares} $r = \cos 2\theta$ (rosácea de quatro pétalas) e $r = \frac{1}{\cos\theta}$ em
$\intoo{-\frac\pi2}{\frac\pi2}$ (reconheça uma reta).
\end{exercise}

\begin{exercise}[$\star\star$]\label{exo:b1:curves:5}
Estude a \omterm{def:b1:curves:polar}{curva polar} $r = 1 + 2\cos\theta$ (uma limaçon): domínio em que $r \geq 0$ contra
$r < 0$ (os pontos marcados com raio negativo ficam no raio oposto),
passagem pela origem, laço interno, esboço.
\end{exercise}

\begin{exercise}[$\star\star$]\label{exo:b1:curves:6}
Identifique a \omterm{def:b1:curves:conics}{cônica} $9x^2 + 25y^2 - 36x - 50y - 164 = 0$: reduza completando quadrados, e dê o
centro, os semieixos, os focos e a excentricidade.
\end{exercise}

\begin{exercise}[$\star\star$]\label{exo:b1:curves:7}
Demonstre que a equação polar de uma \omterm{def:b1:curves:conics}{cônica} com foco na origem é
\[
r = \frac{p}{1 + e\cos\theta}
\]
(diretriz vertical à distância $\frac pe$ à direita do foco). Quais
valores de $\theta$ são permitidos quando $e > 1$?
\end{exercise}

\begin{exercise}[$\star\star$]\label{exo:b1:curves:8}
A partir da propriedade bifocal $MF + MF' = 2a$, deduza a ``construção do
jardineiro'' da elipse, e demonstre que a tangente em $M$ faz ângulos
iguais com $MF$ e $MF'$ \emph{(propriedade de reflexão; derive $\norm{M(t) - F} + \norm{M(t) - F'} = 2a$
e interprete a soma nula de produtos de vetores unitários)}.
\end{exercise}

\begin{exercise}[$\star\star\star$]\label{exo:b1:curves:9}
A \emph{lemniscata de Bernoulli} é a \omterm{def:b1:curves:polar}{curva polar} $r^2 = \cos
2\theta$ (tome $r = \sqrt{\cos2\theta}$
onde estiver definida).
\begin{enumerate}
  \item Dê o seu domínio, as suas simetrias, as tangentes na origem,
        e esboce-a.
  \item Demonstre que ela é o lugar dos pontos $M$ com $MF \cdot MF'
        = \frac12$, onde
        $F, F' = \bigl(\pm\frac{1}{\sqrt2},
        0\bigr)$. \emph{(Calcule $MF^2\,MF'^2$ em \omterm{prop:b1:vspaces:coordinates}{coordenadas} polares.)}
\end{enumerate}
\end{exercise}

\begin{exercise}[$\star\star$]\label{exo:b1:curves:10}
Calcule o \omterm{def:b1:curves:arclength}{comprimento} total da cardioide $r = 1 + \cos\theta$ usando a fórmula polar
da \cref{def:b1:curves:arclength} \emph{(a identidade $1 + \cos\theta = 2\cos^2\frac\theta2$ transforma a raiz quadrada em
$2\abs{\cos\frac\theta2}$)}.
\end{exercise}

\begin{exercise}[$\star\star$]\label{exo:b1:curves:11}
Mostre que a tangente à elipse $(a\cos t, b\sin t)$ no ponto de parâmetro $t$ tem
equação
\[
\frac{x\cos t}{a} + \frac{y\sin t}{b} = 1 ,
\]
e deduza a tangente a $\frac{x^2}{a^2} + \frac{y^2}{b^2} = 1$ num ponto $(x_0, y_0)$ da elipse: $\frac{x\,x_0}{a^2} +
\frac{y\,y_0}{b^2} = 1$ (a regra
do ``desdobramento dos quadrados'').
\end{exercise}

\begin{exercise}[$\star\star\star$]\label{exo:b1:curves:12}
A \emph{espiral logarítmica} é a \omterm{def:b1:curves:polar}{curva polar} $r =
\eu^{k\theta}$ ($k > 0$ fixo,
$\theta \in \R$).
\begin{enumerate}
  \item Mostre que o ângulo $V$ entre o raio e a tangente é
        \emph{constante} ($\tan V = \frac1k$) --- a espiral cruza todo raio sob o
        mesmo ângulo.
  \item Mostre que girar a espiral por um ângulo $c$ a leva na sua
        homotética de razão $\eu^{-kc}$: toda rotação da espiral é uma
        ampliação dela (autossemelhança).
  \item Calcule o \omterm{def:b1:curves:arclength}{comprimento} do arco $\theta \in
        \intoc{-\infty}{\theta_0}$ (como limite dos
        \omterm{def:b1:curves:arclength}{comprimentos} em $\intcc{A}{\theta_0}$, $A \to -\infty$) e observe que ele é finito: uma
        curva que espirala infinitas vezes em torno da origem, de
        \omterm{def:b1:curves:arclength}{comprimento} finito.
\end{enumerate}
\end{exercise}

\section{Problema: a cicloide, rainha das curvas}

\begin{problem}\label{pb:b1:curves:1}
Uma roda de raio $R$ rola sem escorregar ao longo do eixo $x$; o
ponto do aro inicialmente na origem traça a \emph{cicloide}. O século
XVII brigou por causa desta curva --- Galileu pesou recortes de papel
dela, Wren mediu-a, Roberval calculou a sua área, Huygens construiu
relógios sobre ela --- e cada um desses resultados está ao alcance
deste capítulo. Demonstramos os quatro clássicos: a construção da
tangente, o \omterm{def:b1:curves:arclength}{comprimento} $8R$ de Wren, a área $3\pi R^2$, e a propriedade
tautócrona de Huygens.\index{cicloide}\index{tautócrona}

\medskip
\noindent\textbf{Parte I --- O rolamento e a tangente.}
\begin{enumerate}
  \item Depois de a roda ter girado um ângulo $t$, o seu centro fica
        em $\Omega(t) = (Rt, R)$ (rolamento sem escorregamento: distância de contato
        $=$ arco rolado). Mostre que o ponto marcado está em
        \[
        M(t) = \bigl(R(t - \sin t),\; R(1 - \cos t)\bigr).
        \]
  \item Calcule $f'(t)$ e mostre que $\norm{f'(t)} =
        2R\,\abs{\sin\frac t2}$; localize os pontos
        singulares (\omterm{ex:b1:curves:astroid}{cúspides} --- cf.\ o \cref{exo:b1:curves:2}) e o topo de cada
        arco.
  \item Seja $C(t) = (Rt, 0)$ o ponto de contato e $T(t) =
        (Rt, 2R)$ o topo da roda.
        Demonstre que, para $0 < t <
        2\pi$, o vetor $M - C$ é \emph{normal} à
        curva em $M(t)$ e o vetor $T - M$ é \emph{tangente}: para
        traçar a tangente a uma cicloide, ligue o ponto ao topo do
        seu círculo rolante. \emph{(Fatore tudo através de $\sin\frac t2$ e
        $\cos\frac t2$.)}
  \item Interprete a questão 3 cinematicamente: o ponto de contato é
        o \emph{centro instantâneo de rotação}, e a velocidade
        escalar de $M$ é igual à sua distância a $C$ (para velocidade
        angular unitária). Verifique $\norm{M - C} = 2R\abs{\sin\frac t2}$.
  \item Demonstre a relação altura--velocidade
        \[
        \norm{f'(t)}^2 = 2R\,y(t) :
        \]
        numa cicloide percorrida com velocidade angular unitária, a
        velocidade escalar em cada ponto é exatamente a velocidade de
        queda livre para uma descida igual à altura corrente. (Guarde
        isto para a Parte IV.)
\end{enumerate}

\medskip
\noindent\textbf{Parte II --- O teorema de Wren: o arco tem
\omterm{def:b1:curves:arclength}{comprimento} $8R$.}
\begin{enumerate}[resume]
  \item Usando a \cref{def:b1:curves:arclength}, calcule o \omterm{def:b1:curves:arclength}{comprimento} de um arco:
        \[
        L = \int_0^{2\pi} 2R\sin\frac t2\,\dd t = 8R
        \]
        (teorema de Wren, 1658). Quatro diâmetros de roda, e nenhum
        $\pi$ em lugar algum.
  \item Calcule o \omterm{def:b1:curves:arclength}{comprimento de arco} a partir da \omterm{ex:b1:curves:astroid}{cúspide}: $s(t) =
        4R\bigl(1 - \cos\frac t2\bigr)$, e
        verifique $s(2\pi) =
        8R$.
  \item Meça agora o arco a partir do \emph{ápice} $t = \pi$: $\sigma(t) = \abs{4R\cos\frac t2}$.
        Demonstre a relação intrínseca
        \[
        \sigma^2 = 8R\,\bigl(2R - y\bigr) :
        \]
        o quadrado da distância de arco desde o topo é proporcional à
        queda de altura abaixo do topo.
  \item Verificações de coerência: recupere da questão 8 que o meio
        arco do ápice à \omterm{ex:b1:curves:astroid}{cúspide} tem \omterm{def:b1:curves:arclength}{comprimento} $4R$, e compare com o
        cálculo da astroide do \cref{ex:b1:curves:circlelength} --- as duas curvas têm
        \omterm{def:b1:curves:arclength}{comprimentos} algébricos, sem $\pi$; explique o que torna isso
        possível, embora ambas sejam construídas a partir de
        círculos. \emph{(Olhe para a forma de $\norm{f'}$.)}
\end{enumerate}

\medskip
\noindent\textbf{Parte III --- A área de Roberval: $3\pi R^2$.}
\begin{enumerate}[resume]
  \item Justifique que a área entre um arco e o solo é $A = \int_0^{2\pi} y(t)\,x'(t)\,\dd t$ \emph{(a
        substituição $x = x(t)$ em $\int y\,\dd x$, o \cref{thm:b1:integration:parts}; $x$ é crescente)}.
  \item Calcule
        \[
        A = R^2\int_0^{2\pi}(1 - \cos t)^2\,\dd t = 3\pi R^2 :
        \]
        exatamente \emph{três} vezes a área da roda --- a razão que
        Galileu havia adivinhado pesando.
  \item Mesmo método para a astroide $(\cos^3t, \sin^3t)$: mostre que a área
        encerrada é $\frac{3\pi}8$ \emph{(linearize $\sin^4 t\cos^2 t$; só o termo
        constante sobrevive ao longo de um período completo)}.
  \item Confira a fórmula da questão 10 no semicírculo unitário
        superior $(\cos t, \sin t)$, com $t$ indo de $\pi$ a $0$: ela devolve
        $\frac\pi2$?
\end{enumerate}

\medskip
\noindent\textbf{Parte IV --- A tautócrona de Huygens.} Vire o arco:
uma conta desliza sem atrito, sob a gravidade $g$, dentro da tigela
cicloidal
\[
x(t) = R(t + \sin t), \qquad y(t) = R(1 - \cos t)
\qquad (t \in \intcc{-\pi}{\pi}),
\]
cujo ponto mais baixo é a origem ($y$ medido para cima).
\begin{enumerate}[resume]
  \item Calcule a velocidade escalar $\norm{f'(t)} = 2R\cos\frac t2$, o \omterm{def:b1:curves:arclength}{comprimento de arco} a
        partir do fundo $s(t) = 4R\sin\frac t2$, e demonstre a identidade-chave
        \[
        y = \frac{s^2}{8R} .
        \]
  \item A conta solta em repouso a partir do ponto de parâmetro
        $t_0 > 0$ obedece à conservação da energia: se $s(\tau)$
        designa a sua posição de arco no instante $\tau$, então
        $\frac12\bigl(
        \frac{\dd s}{\dd\tau}\bigr)^2 + g\,y = g\,y_0$. Reescreva isto, usando a questão 14, como
        \[
        \Bigl(\frac{\dd s}{\dd\tau}\Bigr)^{\!2}
        = \frac{g}{4R}\,\bigl(s_0^2 - s^2\bigr),
        \qquad s_0 = s(t_0).
        \]
  \item Derive em relação a $\tau$ e obtenha o oscilador harmônico
        \[
        \frac{\dd^2 s}{\dd\tau^2} = -\frac{g}{4R}\,s ;
        \]
        resolva-o com o \cref{thm:b1:diffeq:homogeneous2} e as condições iniciais: $s(\tau) = s_0\cos(\omega\tau)$, $\omega = \sqrt{g/4R}$.
  \item Deduza a \emph{propriedade tautócrona} (Huygens, 1659): o
        tempo para chegar ao fundo,
        \[
        T_{\downarrow} = \frac{\pi}{2\omega}
        = \pi\sqrt{\frac Rg}\,,
        \]
        não depende do ponto de largada --- contas soltas ao mesmo
        tempo de quaisquer duas alturas da tigela chegam juntas.
  \item Verifique por substituição que $s(\tau) =
        s_0\cos\omega\tau$ satisfaz exatamente a
        equação de energia de primeira ordem da questão 15 (e não
        apenas a equação \omterm{def:b1:derivative:def}{derivada}), e explique numa frase por que um
        pêndulo circular é apenas \emph{aproximadamente} isócrono ao
        passo que a cicloide o é exatamente.
  \item Numericamente: qual raio $R$ faz o tempo de descida ser
        exatamente um segundo ($g = 9.81$)? Note quão perto a resposta
        está de um metro, e como o período completo de oscilação
        $2\pi\sqrt{4R/g}$ se compara com a fórmula do pêndulo de pequenas
        oscilações $2\pi\sqrt{\ell/g}$ para $\ell = 4R$.
\end{enumerate}

\medskip
\noindent\textbf{Parte V --- Dividendos, e síntese.}
\begin{enumerate}[resume]
  \item Aplique a regra da tangente da questão 3 em $t =
        \frac\pi2$ (tome
        $R = 1$): calcule $M$, $T$, a direção de $MT$, e confira
        contra $f'(\frac\pi2)$.
  \item (Trocoides) Marque em vez disso um ponto à distância $d$ do
        centro ($d \neq R$): a curva é $x = Rt -
        d\sin t$, $y = R - d\cos t$. Mostre que, para
        $d < R$, a curva é regular em toda parte e não apresenta
        comportamento singular algum ($x' > 0$: ela avança), ao passo
        que, para $d
        > R$, a abscissa $x'$ muda de sinal e a curva faz
        laços --- a roda ferroviária com friso, cujos pontos do aro
        viajam \emph{para trás}.
  \item (Prévia da braquistócrona) Da \omterm{ex:b1:curves:astroid}{cúspide} $(\pi R, 2R)$ da tigela até o
        fundo, compare o tempo de descida pela cicloide $\pi\sqrt{R/g}$ com o
        tempo ao longo da rampa reta que liga os mesmos pontos
        \emph{(aceleração constante $g\sin\alpha$ ao longo da corda)}: mostre
        que a corda leva $\sqrt{\pi^2 + 4}\,\sqrt{R/g} \approx 3.72\sqrt{R/g}$. A curva vence a reta --- ela é, de
        fato, a mais rápida de todas as curvas, um resultado do
        cálculo das variações.
  \item Mostre que, no arco ($0 < t < 2\pi$),
        \[
        \frac{\dd y}{\dd x} = \cot\frac t2,
        \qquad
        \frac{\dd^2 y}{\dd x^2} =
        -\frac{1}{4R\sin^4\frac t2} < 0 :
        \]
        o arco é côncavo, com tangentes verticais exatamente nas
        \omterm{ex:b1:curves:astroid}{cúspides}.
  \item Recupere geometricamente a tangente vertical na \omterm{ex:b1:curves:astroid}{cúspide} do
        \cref{exo:b1:curves:2}: calcule a direção limite da corda $MT$ da questão 3
        quando $t \to
        0^{+}$, sem expansão alguma.
  \item Síntese, em quatro frases: quais três teoremas nomeados este
        problema demonstrou (com os seus números $8R$, $3\pi
        R^2$, $\pi\sqrt{R/g}$ e
        os seus autores); qual único artifício de cálculo (as
        fatorações $\sin\frac t2$, $\cos\frac t2$) alimentou todas as Partes I, II e
        IV; como a relação intrínseca $y = s^2/8R$ converteu geometria numa
        equação diferencial \omterm{def:b1:linmaps:def}{linear}; e em qual capítulo deste livro
        cada Parte se apoiou.
\end{enumerate}
\end{problem}
